% Generated draft: requires mathematical and source review.
\documentclass[11pt]{article}
\usepackage[margin=25mm]{geometry}
\usepackage{fontspec}
\setmainfont{Latin Modern Roman}
\usepackage{amsmath,amssymb,mathtools,unicode-math}
\setmathfont{Latin Modern Math}
\usepackage{enumitem,needspace,xurl,hyperref}
\setlength{\parindent}{0pt}
\setlength{\parskip}{5pt}
\setlength{\emergencystretch}{3em}
\newcommand{\N}{\mathbb{N}}
\newcommand{\Z}{\mathbb{Z}}
\newcommand{\Q}{\mathbb{Q}}
\newcommand{\R}{\mathbb{R}}
\newcommand{\C}{\mathbb{C}}
\newcommand{\F}{\mathbb{F}}
\newcommand{\A}{\mathbb{A}}
\newcommand{\PP}{\mathbb{P}}
\newcommand{\E}{\mathbb{E}}
\newcommand{\eps}{\varepsilon}
\newcommand{\dd}{\,\mathrm{d}}
\begin{document}
\section*{UnsolvedMath ID: 800644}
\textbf{Source problem code:} catalogue entry
\subsection{Definitions and mathematical statement}
The supplied record does not specify the ambient class of spaces or define its starred groupability notation. Consequently, the question below is a conditional formulation requiring those conventions to be recovered from the source.

For a topological space \(X\), the standard unstarred cover notation can be described as follows. An open cover \(\mathcal U\) consists of open subsets of \(X\) whose union is \(X\); under the usual nontrivial-cover convention, \(X\notin\mathcal U\). Write \(\mathcal O(X)\) for these covers. An \(\omega\)-cover is a member \(\mathcal U\in\mathcal O(X)\) such that every finite subset of \(X\) is contained in some member of \(\mathcal U\); write \(\Omega(X)\) for this class. A \(\gamma\)-cover is an infinite member \(\mathcal U\in\mathcal O(X)\) such that each point of \(X\) belongs to all but finitely many members; write \(\Gamma(X)\) for this class. A \(\tau\)-cover is a member \(\mathcal U\in\mathcal O(X)\) such that each point belongs to infinitely many members and, for every \(x,y\in X\), at least one of
\[
 \{U\in\mathcal U:x\in U,\ y\notin U\},\qquad
 \{U\in\mathcal U:y\in U,\ x\notin U\}
\]
is finite. Write \(\mathrm T(X)\) for this class. These descriptions explain standard notation; the supplied record does not confirm its cover conventions, including any countability restriction.

For classes \(\mathcal A(X)\) and \(\mathcal B(X)\) of covers, \(\mathsf S_{\mathrm{fin}}(\mathcal A,\mathcal B)\) means that for every sequence \((\mathcal U_n)_{n\geq1}\) from \(\mathcal A(X)\), there are finite subfamilies \(\mathcal F_n\subseteq\mathcal U_n\) such that
\[
 \bigcup_{n\geq1}\mathcal F_n\in\mathcal B(X).
\]
The property \(\mathsf U_{\mathrm{fin}}(\mathcal A,\mathcal B)\) has the same selection quantifiers restricted to sequences for which no \(\mathcal U_n\) has a finite subcover, and requires instead
\[
 \left\{\bigcup\mathcal F_n:n\geq1\right\}\in\mathcal B(X).
\]

For \(\alpha\in\{\gamma,\omega,\tau\}\), let \(\mathscr G_\alpha(X)\) stand for the source's family denoted \(\mathcal O^{\alpha\text{-*gp*}}\). This is a placeholder for missing source definitions, not a definition of groupability. In particular, the supplied record does not explain the stars or the required grouping conditions.

\textbf{Question (Tsaban, supplied Problem 17).} Let \(\mathscr X\) denote the intended, unspecified class of spaces. With the source's cover and groupability conventions restored, is it true that
\[
 \forall X\in\mathscr X,\qquad
 X\models\mathsf U_{\mathrm{fin}}(\Gamma,\mathrm T)
 \quad\Longleftrightarrow\quad
 X\models\mathsf S_{\mathrm{fin}}(\Omega,\mathscr G_\tau)?
\]
Here \(X\models P\) means that \(X\) has selection property \(P\).

The supplied context recalls the equivalence
\[
 \mathsf U_{\mathrm{fin}}(\Gamma,\mathcal O)
 \Longleftrightarrow
 \mathsf S_{\mathrm{fin}}(\Omega,\mathcal O),
\]
and reports that source reference [coc7] proves
\[
 \mathsf U_{\mathrm{fin}}(\Gamma,\Gamma)
 \Longleftrightarrow
 \mathsf S_{\mathrm{fin}}(\Omega,\mathscr G_\gamma),
\]
while source reference [coc8] proves
\[
 \mathsf U_{\mathrm{fin}}(\Gamma,\Omega)
 \Longleftrightarrow
 \mathsf S_{\mathrm{fin}}(\Omega,\mathscr G_\omega).
\]
These are supplied contextual assertions with unresolved source conventions, not independently verified results here. The registry labels the question open; that label is supplied metadata rather than a verified current literature assessment.
\subsection{Short English statement}
Tsaban's supplied Problem 17 asks whether \(\mathsf U_{\mathrm{fin}}(\Gamma,\mathrm T)\) is equivalent to selecting finitely many sets from each \(\omega\)-cover so that the selected family has the source's \(\tau\)-groupability property. The intended spaces and the starred groupability definition are missing, so the formulation remains conditional and requires human mathematical review.
\subsection{Sources}
\begin{itemize}
\item Local UnsolvedMath registry, permanent ID \(800644\), exact source problem code \texttt{catalogue entry}, collection \emph{AMR Open Problem Lists}; supplied title: \emph{Selection Principles Milestone Problems --- Problem 17}. The record attributes the problem to Tsaban, with proposed year 2003, and labels it open.
\item Tsaban, \emph{Selection Principles in Mathematics: A Milestone of Open Problems} (2003), Problem 17. Supplied source URL: \url{https://arxiv.org/abs/math/0312182}. The record states access on 2026-10-16 and extraction by \texttt{arxiv\_latex}; these are supplied provenance statements, not an independent source visit or verification here.
\item Source-local reference [coc7] is credited in the supplied statement with the equivalence involving \(\mathsf U_{\mathrm{fin}}(\Gamma,\Gamma)\). Its bibliographic entry is not supplied.
\item Source-local reference [coc8] is credited in the supplied statement with the equivalence involving \(\mathsf U_{\mathrm{fin}}(\Gamma,\Omega)\). Its bibliographic entry is not supplied.
\end{itemize}
\end{document}
